%% kapa: exam=1 round=29 year=2018 period=1 form=A subject=civil_law subject_ko=민법 seq=13 number=13
%% source: Q-Net 감정평가사 기출문제 (cst00309) artlSeq=5202808 file=extracted/1008/5202808/2018년 제29회 감정평가사 1차시험 기출문제/2018년 제29회 감정평가사 1차시험 기출문제/1교시 A형.pdf
%% conversion: kordoc 4.24.51 -> tools/gen1.py (auto); review=proofread-sonnet; flags=none
\begin{question}{13}
무권대리와 표현대리에 관한 설명으로 옳은 것은? (다툼이 있으면 판례에 따름)
\begin{choices}
  \item 강행법규에 위반한 무효의 대리행위에 대해서도 표현대리의 법리가 적용될 수 있다.
  \item 무권대리행위의 추인은 본인이 무권대리행위의 상대방뿐만 아니라 무권대리인에 대해서도 할 수 있다.
  \item 상대방의 유권대리 주장에는 표현대리의 성립 역시 포함되므로 법원은 표현대리의 성립 여부까지 판단해야 한다.
  \item 무권대리인이 무권대리행위 후 본인을 단독상속한 경우, 그 무권대리행위가 무효임을 주장하는 것은 신의칙에 반하지 않는다.
  \item 표현대리가 성립하는 경우, 상대방에게 과실이 있으면 과실상계의 법리가 적용된다.
\end{choices}
\end{question}
\answer{2}
